% Generated by task tex-groundtruth from dossiers/gt-flock-ring.md (section 2, A.1 to A.5),
% with the corrections of dossiers/verify-gt-flock-ring.md.
\subsection{The BLAKE2s validity phase (Flock) and ring switching: the transcript as deployed}\label{sec:gen-flock-ring-transcript}

\begin{sloppypar}
Short forms in this section; the number after the colon is the line, at leanVM's pin
\code{a386121f}.
\begin{itemize}
\item Specification, under \file{doc/leanvm/body/}: \code{c-flock:} is
  \file{c-flock-protocol.tex}, \code{a-ring:} \file{a-ring-switching.tex}, \code{03:}
  \file{03-proving-primitives.tex}, \code{04:} \file{04-committing-the-witness.tex}, \code{07:}
  \file{07-instruction-tables.tex}, \code{08:} \file{08-end-to-end-protocol.tex}.
\item Rust, under \file{crates/flock/src/}: \code{hash.rs}, \code{zerocheck.rs},
  \code{lincheck.rs}, \code{reduction_tests.rs}, \code{flock/src/lib.rs}; under
  \file{crates/flock/src/zerocheck/}: \code{univariate_skip_optimized.rs},
  \code{zerocheck/multilinear.rs}.
\item Rust, under \file{crates/lean_vm/src/}: \code{cpu/mod.rs}, \code{cpu/layout.rs},
  \code{hash_flock.rs}, \code{tables.rs}; under \file{crates/pcs/src/}: \code{pack.rs},
  \code{stack_open.rs}, \code{ring_switch.rs}, \code{tensor_algebra.rs};
  \code{transcript.rs} is \file{crates/fiat_shamir/src/transcript.rs}.
\item Python: \code{verifier.py} is \file{python-verifier/verifier.py}.
\end{itemize}
\end{sloppypar}

\paragraph{Notation.} $k_{\mathrm{batch}}$ is the log of the number of BLAKE2s compressions (the
log-height $τ$ of the \code{BLAKE2S} table), $k_{\mathrm{bool}} = 14 + k_{\mathrm{batch}}$ the
number of variables of the bit witness $z$, $n_{\mathrm{flock}} = k_{\mathrm{bool}} − 6 = 8 +
k_{\mathrm{batch}}$ the number of variables of the packed witness $\qflock$.

\subsubsection{What is committed}\label{sec:gen-flock-ring-committed}

Nothing in this phase. The bit witness $z$ was committed in the commit phase, packed 64 bits to a
$\K$ element, as the region $\qflock$ of the one stack $q$.

{\small
\setlength{\tabcolsep}{4pt}
\begin{longtable}{@{}>{\raggedright\arraybackslash}p{0.21\linewidth}>{\raggedright\arraybackslash}p{0.26\linewidth}>{\raggedright\arraybackslash}p{0.29\linewidth}>{\raggedright\arraybackslash}p{0.18\linewidth}@{}}
\caption{What is committed for the Flock phase.}\label{tab:gen-flock-ring-committed}\\
\toprule
& Specification & Rust & Python \\
\midrule
\endfirsthead
\toprule
& Specification & Rust & Python \\
\midrule
\endhead
\bottomrule
\endfoot
$z$ is packed into $\qflock$, a region of the stack
& \code{c-flock:30} \enquote{The resulting multilinear $\qflock$ \dots\ is committed as one region
  of the stack}; \code{08:60}
& \code{hash_flock.rs:3-5} \enquote{committed as a column in leanVM's ONE stacked \code{F64}
  witness \dots, with no separate flock commitment}; \code{cpu/layout.rs:25} \code{QFLOCK = 5}
& \code{verifier.py:633}, \code{:879-880} \\ \addlinespace[2pt]
Packing: bit $i$ of word $u$ is $z(i, u)$, coefficient of $x^i$
& \code{a-ring:13-16} $\qflock(u) = Σ_i Q_i(u) x^i$
& \code{pack.rs:1-2} \enquote{least significant bit first}; \code{hash_flock.rs:187-197}
& (verifier only) \\ \addlinespace[2pt]
No other oracle, no other commitment
& \code{c-flock:264} \enquote{the PCS proves the claims opened on $\qflock$}
& \code{flock/src/lib.rs:6-13}; \code{cpu/mod.rs:564-571}
& \code{verifier.py:1404-1413} \\
\end{longtable}
}

The vectors $a = Az$ and $b = Bz$ are never committed and never opened: only $z$ is
(\code{c-flock:30}; in the Rust the only commitment is \code{pcs::commit(&mut ps, &w.q, …)},
\code{cpu/mod.rs:560-562}, and $a$, $b$ are not columns of the layout). (The citation check
corrected the support of this sentence: the dossier cited \code{c-flock:18}, which only defines
$a = Az$, $b = Bz$, $c = Cz = z$, and \code{hash.rs:592-593}, \enquote{No c buffer. Since
\code{C = I}, \code{c == z}}, which supports only $c = z$.)

\subsubsection{The messages, in order}\label{sec:gen-flock-ring-messages}

\code{P →} is a prover message, \code{V →} a verifier challenge, \code{V:} a verifier
computation. All values are in $\E$. Counts are per proof.

\begin{landscape}
{\small
\setlength{\tabcolsep}{4pt}
\begin{longtable}{@{}>{\raggedright\arraybackslash}p{0.03\linewidth}>{\raggedright\arraybackslash}p{0.33\linewidth}>{\raggedright\arraybackslash}p{0.19\linewidth}>{\raggedright\arraybackslash}p{0.25\linewidth}>{\raggedright\arraybackslash}p{0.15\linewidth}@{}}
\caption{The Flock phase and ring switching, message by message, as deployed.}\label{tab:gen-flock-ring-messages}\\
\toprule
\# & Step & Specification & Rust (verifier unless said) & Python \\
\midrule
\endfirsthead
\toprule
\# & Step & Specification & Rust (verifier unless said) & Python \\
\midrule
\endhead
\bottomrule
\endfoot
1 & \code{V:} seven fixed coordinates $r_0..r_6$ = $φ_8(\mathtt{0xf7}), φ_8(\mathtt{0x53}),
  φ_8(\mathtt{0xb5}), g_0^{2^j}/(1+g_0^{2^j})$, $j<4$
  & \code{c-flock:54} (value of $g_0$ not given: \enquote{where $g_0∈\E$ is public})
  & \code{zerocheck.rs:51-58}; \code{univariate_skip_optimized.rs:65, 79-101}; $g_0$ at
  \code{:104-106}
  & \code{verifier.py:1095-1100} \\ \addlinespace[2pt]
2 & \code{V →} $n_{\mathrm{flock}} − 7 = k_{\mathrm{batch}} + 1$ coordinates $r_7..$, before any
  Flock message
  & \code{c-flock:55}, \code{:173}
  & \code{zerocheck.rs:362} \code{equality_tail(m, |n| vs.sample_vec(n))}
  & \code{verifier.py:1139} \\ \addlinespace[2pt]
3 & \code{P →} $P$ on the coset $H' = φ_8(64..127)$: 64 values
  & \code{c-flock:62} \enquote{The prover sends $P|_{H'}$, $64$ values}
  & \code{zerocheck.rs:365} (prover \code{:185-187})
  & \code{verifier.py:1142} \\ \addlinespace[2pt]
4 & \code{V →} $z_{\mathrm{skip}}$
  & \code{c-flock:64}
  & \code{zerocheck.rs:366}
  & \code{verifier.py:1143} \\ \addlinespace[2pt]
5 & \code{V:} $v_P$, the value at $z_{\mathrm{skip}}$ of the interpolant of the 64 received values
  and \textbf{64 assumed zeros} on $H$; no check
  & \code{c-flock:62-67}
  & \code{zerocheck.rs:378}; \code{zerocheck/multilinear.rs:129-138}
  & \code{verifier.py:1144} \\ \addlinespace[2pt]
6 & $n_{\mathrm{flock}}$ rounds. \code{P →} two coefficients $c_1, c_2$ of the quadratic cofactor
  $G_i$; \code{V:} $c_0 := \mathit{claim} + r_i·(c_1 + c_2)$; \code{V →} $χ_i$; \code{V:}
  $\mathit{claim} := G_i(χ_i)$. Low variable first. No check
  & \code{c-flock:69} (\enquote{with Gruen's eq-factor optimization}); \code{03:110}
  & \code{zerocheck.rs:397-405}; \code{transcript.rs:289-309} (prover \code{zerocheck.rs:116-123})
  & \code{verifier.py:1147}, \code{:406-413}, \code{:420-427} \\ \addlinespace[2pt]
7 & \code{P →} $v_a, v_b$; \code{V:} $v_c := R_{\mathrm{zc}} + v_a·v_b$ (never sent, never checked)
  & \code{c-flock:69-73} \enquote{is what \emph{defines} $v_c$: both sides solve it rather than
  transmit it}
  & \code{zerocheck.rs:421-423}
  & \code{verifier.py:1148-1149} \\ \addlinespace[2pt]
8 & \code{V →} $α_{\mathrm{lc}}$; \code{V:} target $v_a + α v_b + α² v_c + α³$
  & \code{c-flock:112-118}
  & \code{lincheck.rs:1179}, \code{:1199-1203}
  & \code{verifier.py:1155}, \code{:1162} \\ \addlinespace[2pt]
9 & 8 rounds. \code{P →} two coefficients $c_0, c_2$ of the quadratic; \code{V:} $c_1 :=
  \mathit{claim} + c_2$; \code{V →} challenge; high variable first, so the point
  $χ'_{\mathrm{in}}$ is the challenges reversed
  & \code{c-flock:122}
  & \code{lincheck.rs:1205-1211}, \code{:1218-1219} (prover \code{:1085-1107})
  & \code{verifier.py:1163}, \code{:1169} \\ \addlinespace[2pt]
10 & \code{P →} $s_0..s_{63}$, claimed $s_i = \mle{z}(i, χ'_{\mathrm{in}}, χ_{\mathrm{out}})$
  & \code{c-flock:124}
  & \code{lincheck.rs:1214} (prover \code{:1113-1115})
  & \code{verifier.py:1168} \\ \addlinespace[2pt]
11 & \code{V:} \textbf{check} the terminal identity $R_{\mathrm{lc}} = e_{\mathrm{row}}^T (A_0 + α
  B_0) w_{\mathrm{col}} + α²·\eq(χ_{\mathrm{in}}, χ'_{\mathrm{in}})·Σ_i ϖ_i(z_{\mathrm{skip}})
  s_i + α³·\eq(χ'_{\mathrm{in}}, 8)·s_0$; reject otherwise
  & \code{c-flock:126-133}, \code{:150}
  & \code{lincheck.rs:1233-1256}, error \code{ConsistencyFailed}, surfaced as
  \code{CpuError::Blake2s} at \code{cpu/mod.rs:765}; the matrix term is \code{hash.rs:361-389}
  (\code{bilinear_walk})
  & \code{verifier.py:1170-1176} \code{require(terminal == r_lc, …)}; matrix term
  \code{:1180-1301} \\ \addlinespace[2pt]
12 & \code{V →} six challenges $f_0..f_5$ (ring switching)
  & \code{a-ring:104}; \code{08:91}
  & \code{stack_open.rs:513}; \code{ring_switch.rs:160-162}
  & \code{verifier.py:1343} \\ \addlinespace[2pt]
13 & \code{V:} target $T = Σ_i x^i Φ(s_i)$; the weight is $W(u) = Φ(\eq(r, u))$ on the
  $\qflock$ region, $r = (χ'_{\mathrm{in}}, χ_{\mathrm{out}})$. No message, no check
  & \code{a-ring:44-46}, \code{:104-110}
  & \code{stack_open.rs:521-524}; \code{ring_switch.rs:555-557} (through the transposed columns,
  \code{tensor_algebra.rs:42-71}, and \code{build_coordinate_weights},
  \code{ring_switch.rs:133-144})
  & \code{verifier.py:1345-1347} \\ \addlinespace[2pt]
14 & (opening) \code{V →} $λ$; the ring-switched claim takes $λ^0$, the pooled claims the next
  powers
  & \code{08:98}
  & \code{stack_open.rs:518-527}
  & \code{verifier.py:1361}, \code{:1413} \\ \addlinespace[2pt]
15 & (opening, last step) \code{V:} evaluates $\mle{W}$ at the final point $r'$
  & \code{a-ring:150-154}: $Σ_k c_k ∏_n (1 + r_n^{2^k} + r'_n)$
  & \code{stack_open.rs:530-546}; \code{ring_switch.rs:597-617} (\code{eval_rs_eq}, the
  tensor-algebra algorithm of Diamond--Posen)
  & \code{verifier.py:1324-1335}, \code{:1412} (the closed form) \\
\end{longtable}
}
\end{landscape}

Counts, checked against \code{reduction_tests.rs:166-168} and by the probe of the dossier's
section 9.3 (the citation check re-derived both from the code):

\begin{center}\small
\begin{tabular}{@{}>{\raggedright\arraybackslash}p{0.36\linewidth}>{\raggedright\arraybackslash}p{0.58\linewidth}@{}}
\toprule
What & Count \\
\midrule
Scalars the reduction (steps 3 to 10) puts on the stream
  & $64 + 2·n_{\mathrm{flock}} + 2 + 16 + 64 = 162 + 2·k_{\mathrm{batch}}$ (168 at
  $k_{\mathrm{batch}} = 3$) \\
Challenges drawn by steps 2 to 12
  & $(k_{\mathrm{batch}} + 1) + 1 + n_{\mathrm{flock}} + 1 + 8 + 6 = 2·k_{\mathrm{batch}} + 25$ \\
\bottomrule
\end{tabular}
\end{center}

Steps 12 and 13 are \enquote{BLAKE2s validity, step 2} in the specification (\code{08:91}) and
the first lines of the opening function in the Rust (\code{open_batch_mixed_whir_stacked},
\code{verify_opening_batch_mixed_whir_stacked}). The transcript position is the same.

\subsubsection{Values derived, never transmitted}\label{sec:gen-flock-ring-derived}

$v_P$ (step 5); $c_0$ of each of the $n_{\mathrm{flock}}$ zerocheck rounds (step 6); $v_c$
(step 7); the lincheck target (step 8); $c_1$ of each of the 8 lincheck rounds (step 9); $T$
(step 13). The $s_i$ are sent once, by the lincheck, and reused by ring switching
(\code{stack_open.rs:113-116} \enquote{this layer reads nothing off the stream}).

\subsubsection{The eighteen limb claims}\label{sec:gen-flock-ring-limbs}

They are produced by the table sumcheck, routed to the opening, and never seen by Flock.

{\small
\setlength{\tabcolsep}{4pt}
\begin{longtable}{@{}>{\raggedright\arraybackslash}p{0.19\linewidth}>{\raggedright\arraybackslash}p{0.24\linewidth}>{\raggedright\arraybackslash}p{0.31\linewidth}>{\raggedright\arraybackslash}p{0.20\linewidth}@{}}
\caption{The eighteen limb claims, from the table sumcheck to the opening.}\label{tab:gen-flock-ring-limbs}\\
\toprule
& Specification & Rust & Python \\
\midrule
\endfirsthead
\toprule
& Specification & Rust & Python \\
\midrule
\endhead
\bottomrule
\endfoot
The table sumcheck sends one value per column of every table, the eighteen limb columns included
& \code{08:77} \enquote{The prover sends one value per column of every table}
& \code{cpu/mod.rs:428-441} (\code{for c in 0..table.n_committed_columns()}, 37 for
  \code{BLAKE2S}, \code{tables.rs:866-868}, \code{:862})
& \code{verifier.py:620} (\code{table.width}), \code{:827-834} \\ \addlinespace[2pt]
A limb claim at $z$ becomes a claim on $\qflock$ with the low 8 coordinates frozen to the slot
& \code{08:77} \enquote{their claims route there}; \code{07:122}. The selector is \textbf{not
  described} (§4.1 has aligned blocks only, \code{04:10-18})
& \code{cpu/mod.rs:790-814}, which builds a \code{pcs::SlotClaim::Strided} (quoted below the
  table); \code{stack_open.rs:84-97}, \code{:305-323}; \code{hash_flock.rs:223} (\code{= 14 − 6
  = 8})
& \code{verifier.py:884-894} \code{Placement(kappa, offsets[QFLOCK] + limbs[column],
  QFLOCK_SLOT_BITS)}; \code{:295-302} \\ \addlinespace[2pt]
Slot of each limb
& not given
& \code{hash_flock.rs:87-115}: \code{cv} 0..3, \code{out} 4..7, message 10..17, metadata 18, 19
& \code{verifier.py:847-850} (same) \\ \addlinespace[2pt]
The Flock verifier takes no claim
& \code{c-flock:170-179} (the summary has no input claim)
& \code{cpu/mod.rs:765} \code{verify_reduction(n_blocks, &mut vs)}; \code{hash_flock.rs:282-284}
& \code{verifier.py:1405}, called with the size only (quoted below the table); the signature
  \code{verify_flock(log_n: int, transcript: Transcript)} is at \code{:1304} \\ \addlinespace[2pt]
The limb claims enter the opening batch with every other pooled claim
& \code{08:97-99}
& \code{cpu/mod.rs:756, 768} (\code{slots}), \code{stack_open.rs:525-527}
& \code{verifier.py:1395, 1413} \\
\end{longtable}
}

The Rust's strided claim, and the Python's call of the Flock verifier:

\begin{srccode}
return pcs::SlotClaim::Strided {
    offset: l.placements[QFLOCK].offset,
    slot,
    stride_log: crate::hash_flock::SLOT_STRIDE_LOG,
    point: c.point.clone(),
    value: c.value,
};
\end{srccode}
\src{crates/lean\_vm/src/cpu/mod.rs:799-805}

\begin{srccode}
flock_point, flock_s = verify_flock(BLAKE2S_R1CS_LOG_SIZE + layout.table_log_heights[OP_BLAKE2S], transcript)
\end{srccode}
\src{python-verifier/verifier.py:1405}

\begin{sloppypar}
(The citation check corrected two quotations of the dossier: \code{cpu/mod.rs:790-814} was
quoted in abbreviated form as \code{SlotClaim::Strided { offset, slot, stride_log:
SLOT_STRIDE_LOG, point, value }}, and \code{verifier.py:1405} as \code{verify_flock(log_n,
transcript)}, which is the signature at \code{:1304}.)
\end{sloppypar}

\subsubsection{What leaves the phase}\label{sec:gen-flock-ring-output}

One family of 64 slice claims at one point, turned by ring switching into \textbf{one weighted
claim} $⟨W, q⟩ = T$ whose weight is supported on the $\qflock$ region (lifted to the stack by the
region's selector: \code{stack_open.rs:41-48}, \code{verifier.py:1412}), plus the pool unchanged.
